\section{Gluing a left state at a cut of the chain}
\label{sec:chain_gluing}

Fix the cut after $k$ sites of a chain. Splitting off the last site left of
the cut, a configuration has coordinates $((x,y),z)$, with $x$ on the first
$k-1$ sites, $y$ on site $k-1$ and $z$ on the sites right of the cut. This is
the setting of the boundary contractions and the gluing lemma of Landau,
Vazirani and Vidick~\cite[lines 409--447]{Landau2015Polynomial}, proved in the
companion quantum-channel
volume~\cite[``Gluing'', ``The leading eigenvector of a near-optimal left state'',
    ``Three-way cut coordinates'']{QICLean2026Blueprint}.

\begin{definition}[Cut coordinates at a cut of the chain]
    \label{def:chain_gluing_cut}
    \lean{TNLean.MPS.AreaLaw.leftBlock_sub_one_subset, TNLean.MPS.AreaLaw.CutLeft,
        TNLean.MPS.AreaLaw.CutSite, TNLean.MPS.AreaLaw.CutRight,
        TNLean.MPS.AreaLaw.gluingCut, TNLean.MPS.AreaLaw.leftCutOperator,
        TNLean.MPS.AreaLaw.rightCutOperator}
    \leanok
    The coordinates $((x,y),z)$ of a configuration of the chain at the cut
    after $k$ sites. For an operator $H_L$ supported on the first $k$ sites,
    $H_L'$ is its restriction to them; for an operator $K$ supported off the
    first $k-1$ sites, $K'$ is its restriction to site $k-1$ and the sites right
    of the cut.
\end{definition}

\begin{theorem}[The Hamiltonian at a cut]
    \label{thm:chain_gluing_split}
    \lean{TNLean.MPS.AreaLaw.reindex_gluingCut_eq}
    \leanok
    \uses{def:chain_gluing_cut}
    If $H_L$ is supported on the first $k$ sites and $K$ off the first $k-1$
    sites, then in the cut coordinates $H_L+K=H_L'\otimes1+1\otimes K'$.
\end{theorem}

\begin{proof}
    \leanok
    An operator supported on a set of sites acts as the identity on the
    configuration of the other sites. Thus, in the cut coordinates, with
    the reassociation $((x,y),z)\leftrightarrow(x,(y,z))$ for $K$,
    \begin{align}
        H_L & =H_L'\otimes1, \\
        K   & =1\otimes K'.
    \end{align}
    Adding gives the asserted decomposition of $H_L+K$.
\end{proof}

\begin{theorem}[The leading eigenvector at a cut of the chain]
    \label{thm:chain_leading_eigenvector}
    \lean{TNLean.MPS.AreaLaw.exists_leading_eigenvector_chain}
    \leanok
    \uses{def:chain_gluing_cut}
    Let $H=H_L+K$ with $H_L$ supported on the first $k$ sites and $K$ off the
    first $k-1$ sites have a gapped ground state $\Gamma$ with energy $E_0$ and
    gap $\Delta$. Let $W$ be an isometry from a bond $\mathbb C^m$ into the
    sites right of the cut, and let $w=(1\otimes W)x$ in the cut coordinates
    have energy at most $E_0+e$. Let $\sigma$ be a density matrix on the first
    $k$ sites and the bond with
    $\operatorname{Re}\operatorname{tr}(\sigma(H_L'\otimes1))\le
        \operatorname{Re}\operatorname{tr}(\lvert x\rangle\langle x\rvert(H_L'\otimes1))$
    and
    $\lVert\operatorname{tr}_{[1,k-1]}\sigma-\operatorname{tr}_{[1,k-1]}
        \lvert x\rangle\langle x\rvert\rVert_1\le\tau$. If
    $4(e+\lVert K\rVert\tau)<\Delta$, then $\sigma$ has a unit eigenvector $u$
    with eigenvalue larger than $\tfrac12$ such that $(1\otimes W)u$, read on
    the chain, has energy at most $E_0+2(e+\lVert K\rVert\tau)$ and overlap at
    least $1-2(e+\lVert K\rVert\tau)/\Delta$ with $\Gamma$.

    This is Lemma~l:leading of~\cite[lines 493--505]{Landau2015Polynomial} on
    the chain, with $H_L=H_1+\cdots+H_{i-1}$ and $K=H_i+H_R$; the energy excess
    $e$ of the trimmed state of Lemma~l:within is corrected
    in~\cite{gap:lvv13_trimming_energy}.
\end{theorem}

\begin{proof}
    \leanok
    \uses{thm:chain_gluing_split}
    In the cut coordinates the Hamiltonian, the ground state and $w$ satisfy
    the hypotheses of the leading-eigenvector theorem of the companion volume,
    with $\lVert K'\rVert\le\lVert K\rVert$ since $K'$ is a corner of
    $1\otimes K'$. Let $E$ be the cut-coordinate bijection and $U$ the
    unitary map $U\psi=\psi\circ E^{-1}$. For chain vectors $\phi,\psi$ and
    an operator $M$,
    \begin{align}
        \langle U\phi,U\psi\rangle
         & =\langle\phi,\psi\rangle,
        \label{eq:chain_gluing_inner} \\
        \langle U\phi,(UMU^*)U\psi\rangle
         & =\langle\phi,M\psi\rangle.
        \label{eq:chain_gluing_energy}
    \end{align}
    Applying~(\ref{eq:chain_gluing_inner}) and~(\ref{eq:chain_gluing_energy})
    to the ground state and the eigenvector
    extension carries the energy and overlap bounds back to the chain.
\end{proof}
